\section{对手给数据时该比什么}
\label{sec:17-Constrained-Guarantees/04-Online-Learning-and-Regret/01}

一套反欺诈规则每周重训一次。离线指标一直漂亮：留出集上的召回率稳定在 \(0.86\) 上下，上线前的复核也没查出泄漏。可线上拦截率从第一周的 \(0.83\) 一路滑到第六周的 \(0.51\)，中间没有任何一次发版改动过模型。

排查的第一反应是分布漂移。但漂移通常有形状——某个渠道的流量涨了，某个字段的口径变了，某个季节性因素起来了。这里查不到这类东西。真正在变的是被拦截样本的构造方式：每次上线一批新规则，一周之内那批模式就消失，取而代之的是一批只在被拦的那几个字段上做了微调的新模式。改动精确地落在起作用的特征上，其余部分原封不动。

这不是自然界的变化，是有人在读你的拦截结果并据此调整。你堵住一个口子，他们换一个。

一旦意识到这一点，``模型在测试集上的误差是多少''这句话就从根上失效了。测试集是过去的样本，而对手是活的：他下一周的行为依赖于你这一周的部署。留出集再干净，它也只能告诉你``如果世界还是上周那个样子''会怎样，而世界不再是上周那个样子，恰恰是因为你部署了模型。\hyperref[sec:18-Synthesis/01-System-Lifecycle/06]{``上线之后系统开始改变它所预测的世界''}讲的是这件事的一般形态，反欺诈只是其中最尖锐的一种——那里的反馈回路是无意的，这里的反馈回路背后坐着一个想赢的人。

\subsection{i.\,i.\,d.\ 是第一个塌掉的东西}

统计学习理论给出的所有保证，都挂在同一个前提上：训练样本与测试样本独立同分布地来自同一个未知分布 \(D\)。\hyperref[sec:13-Machine-Learning/10-Learning-Theory/01]{``训练误差为何总是乐观''}用一百万枚硬币说明了为什么需要一致收敛，\hyperref[sec:13-Machine-Learning/10-Learning-Theory/02]{``从联合界到一致收敛''}把它做成了对整个假设类都成立的偏差控制。但这些推导的第一行永远是同一句：设 \(S=\{z_1,\dots,z_n\}\) 独立同分布地抽自 \(D\)。

对抗环境里，这一行就不成立。样本之间不独立，因为第 \(t+1\) 个样本是对手看过前 \(t\) 轮结果之后造出来的；也不同分布，因为对手的策略在变。这两条一塌，\hyperref[sec:05-Probability/08-Limit-Theorems/02]{``大数定律与长期平均''}里那句``样本均值向期望靠拢''就失去依据，用留出集上的平均误差去估计真风险这一步不再合法。更要命的是，根本不存在一个固定的 \(D\) 让我们去讨论``真风险''。风险是关于某个分布的期望，分布都没有了，期望自然无从谈起。

这不是把界放宽一点就能补救的事。\hyperref[sec:13-Machine-Learning/10-Learning-Theory/04]{``泛化界的解释力与边界''}提醒过，泛化界的数值本来就宽松，它的价值在于指出哪些因素影响泛化。而在这里，连``影响什么''都失去了指代对象：没有 \(D\)，就没有 \(R(h)=\E_{z\sim D}\ell(h,z)\) 这个被影响的量。\hyperref[sec:13-Machine-Learning/01-Learning-Problems-and-Evaluation/04]{``数据划分、泄漏与分布偏移''}里那套划分技术管的是评估流程不要作弊，它假定世界本身是被动的，而这里世界会回话。

所以问题不是``怎样把泛化界修好''，是``换一个什么承诺，才能在完全不知道数据怎么来的情况下仍然兑现''。

\subsection{把承诺换成事后对比}

在线学习把问题重写成一个逐轮进行的协议。共 \(T\) 轮，第 \(t\) 轮里我们先从决策集 \(\mathcal X\) 中选出 \(x_t\)，然后对手给出一个损失函数 \(\ell_t\)，我们承受 \(\ell_t(x_t)\)。对手可以在看到 \(x_1,\dots,x_{t-1}\) 之后再决定 \(\ell_t\)；不做任何关于 \(\ell_t\) 序列的概率假设，它可以是完全恶意的。

要在这样的设定下承诺什么？``总损失小于某个绝对值''显然做不到——对手可以让每个可选动作都很糟。可行的办法是把承诺改成相对的：

\[
R_T=\sum_{t=1}^{T}\ell_t(x_t)\;-\;\min_{x\in\mathcal X}\sum_{t=1}^{T}\ell_t(x).
\]

第一项是我们实际付出的累计损失。第二项是一个事后诸葛亮：等 \(T\) 轮全部结束、所有 \(\ell_t\) 都摊在桌上之后，回过头去挑一个自始至终不变的决策 \(x\)，让它的累计损失最小。两者之差叫\emph{遗憾}——回头看，早知道就该一直用那个 \(x\)，少受的那部分损失。

这个定义的巧妙之处在于它对世界一无所求。它不假设分布存在，不假设样本独立，不假设对手弱。无论 \(\ell_1,\dots,\ell_T\) 是怎么来的，公式两边用的是同一串损失函数，同一个对手的同一套出招。所以``\(R_T\) 增长得比 \(T\) 慢''是一个即使对手完全恶意也能兑现的承诺，而``测试误差不超过多少''不是。

\begin{center}
\begin{tikzpicture}[x=1cm,y=1cm,font=\small,>=stealth]
  \draw[->,black!60] (0,0) -- (7.4,0) node[right,font=\footnotesize] {\(t\)};
  \draw[->,black!60] (0,0) -- (0,4.4) node[above,font=\footnotesize] {累计损失};
  \draw[blue!65!black,thick]
    (0,0) -- (1,0.69) -- (2,1.27) -- (3,1.83) -- (4,2.38) -- (5,2.92) -- (6,3.47) -- (7,4.00);
  \draw[orange!75!black,thick]
    (0,0) -- (7,3.50);
  \draw[<->,black!70] (7,3.50) -- (7,4.00);
  \node[anchor=west,font=\footnotesize,black!75] at (7.1,3.75) {\(R_T\)};
  \draw[blue!65!black,thick] (4.2,1.45) -- (4.65,1.45);
  \node[anchor=west,font=\footnotesize] at (4.72,1.45) {算法实际付出};
  \draw[orange!75!black,thick] (4.2,0.85) -- (4.65,0.85);
  \node[anchor=west,font=\footnotesize] at (4.72,0.85) {事后最优固定决策};
  \foreach \x in {1,3,5,7} \draw[black!45] (\x,0) -- (\x,-0.1) node[below,font=\scriptsize] {\x};

  \begin{scope}[shift={(9.0,0.5)}]
    \draw[->,black!60] (0,0) -- (3.3,0) node[right,font=\footnotesize] {\(T\)};
    \draw[->,black!60] (0,0) -- (0,2.5) node[above,font=\footnotesize] {\(R_T/T\)};
    \draw[blue!65!black,thick,domain=0.30:3.1,samples=60]
      plot ({\x},{0.62/sqrt(\x)});
    \draw[black!35,dashed] (0,0) -- (3.1,0);
    \node[anchor=west,align=left,font=\scriptsize,black!70] at (0.9,1.55)
      {平均遗憾趋于零\\ 不代表损失趋于零};
  \end{scope}
\end{tikzpicture}

\emph{左图里两条曲线之间的竖直间隙就是遗憾；右图说明我们真正要的是这个间隙除以 \(T\) 之后归零，也就是在平均意义上追平那条橙线。}
\end{center}

图里橙线是直的，因为固定决策每轮承受的损失取自同一个 \(x\)；蓝线在早期爬得快，那段是算法还没认清局势时多付的部分。两条线之间的间隙一直在长，但长得越来越慢——这正是 \(R_T=o(T)\) 的图像含义。称一个算法\emph{无遗憾}，指的就是对任意对手序列都有 \(R_T/T\to0\)。

\subsection{这个承诺弱在哪里}

\(R_T/T\to0\) 读作``长期平均下来，我们不比那条橙线差''。这句话里每一个限定词都在削弱它，必须逐个说清，否则很容易把它当成比实际更强的东西。

橙线本身可能很低。遗憾是相对量，它只保证我们追平事后最优的\emph{固定}决策，完全没说那个决策好不好。若反欺诈的整个规则空间里，最好的那条固定规则也只拦得住四成欺诈，那么一个无遗憾算法的长期表现就是四成上下。它兑现了承诺，业务上依然不合格。把这件事说反过来更清楚：无遗憾算法保证``不比某个基准差''，而基准是我们自己选的，选得太弱，保证就没有内容。

比较类只含固定决策。橙线是一条直线，因为 \(\min\) 是在单个 \(x\) 上取的。比较类的大小在这里扮演的角色，与假设类容量在泛化界里扮演的角色是同一个——\hyperref[sec:13-Machine-Learning/10-Learning-Theory/03]{``VC 维把无限假设类变有限''}衡量的是后者，遗憾界里那个 \(\ln N\) 或 \(\ln k\) 衡量的是前者。若真实世界前半段适合规则 A、后半段适合规则 B，那么``事后最优固定决策''这个基准两段都不合身，追平它并不意味着我们跟上了变化。想把基准换成``最优的 \(k\) 段分段固定决策''是可以做的，遗憾界里会多出一个与 \(k\) 有关的因子，为此让出的是常数——基准每强一档，能保证的界就松一档。这条``基准强度与可达界成反比''的规律贯穿整个在线学习。

只有平均意义。\(R_T/T\to0\) 不排除某一轮、某一段时间里表现极差。它是把整段历史摊平之后的陈述，对任何单轮不作保证。在损失可以致命的场合，摊平之后的承诺不够用。

保证是对\emph{最坏}对手的。这既是它的强处也是它的限制：因为要对付所有可能的对手序列，界里不含任何``实际情况通常没那么坏''的红利。真实数据往往比最坏情况温和得多，此时无遗憾算法会显得保守。近年的自适应遗憾界试图让界随实际难度收缩，但它们仍以最坏情况为上限。

把这四条并排放，遗憾的定位就清楚了：它不是一个更好的泛化界，而是一个换了参照系的承诺。泛化界断言的是``我们与真风险的差距''，参照系在世界那边；遗憾断言的是``我们与一个事后基准的差距''，参照系被搬到了我们自己这边。搬过来之后，关于世界的假设就都不需要了，能说的话也随之变弱。这是一次明确的交换，不是免费的改进。

那么剩下的问题只有一个：这个看起来很谦逊的承诺，真的做得到吗？对手可以任意出招，我们却要保证累计损失追平一个事后才知道的基准，直觉上这像是在要求预知未来。下一节给出两个具体算法，说明不需要预知未来，也说明各自要付出什么条件才行。
